\chapter{Conclusiones y trabajo futuro}
\label{chap:conclusiones}

\drop{D}{esde} el punto de vista técnico, \dataqual aporta al panorama de
herramientas de calidad de datos una solución \emph{full-stack} autoalojada y
sin necesidad de programar que reúne métricas específicas para \ac{ML},
\emph{fingerprinting} determinista, \emph{Quality Gates} y protección nativa de
\ac{PII}, una combinación ausente en las herramientas de código abierto
analizadas en el estado del arte (sección~\ref{sec:estado-arte}) y cuyo detalle
se recoge en la sección~\ref{sec:fortalezas}.

Desde el punto de vista académico, el proyecto materializa la convergencia entre
la ingeniería del software moderna (arquitecturas de microservicios, patrones de
diseño como \emph{plugin} y factoría) y el dominio emergente de la calidad de
datos para \ac{IA}, proporcionando una implementación de referencia alineada
con los estándares \ac{ISO}/\ac{IEC}~25012 e \ac{ISO}/\ac{IEC}~5259 que puede servir como
base para investigación futura sobre métricas de calidad específicas para
\ac{ML}.



%% =========================================================================
\section{Revisión de objetivos}
\label{sec:revision-objetivos}
%% =========================================================================

El presente \ac{TFG} estableció cinco objetivos específicos (OE1--OE5) y
un conjunto de requisitos no funcionales.  La verificación detallada del
grado de cumplimiento, con la evidencia técnica asociada a cada uno, se
recoge en los Cuadros~\ref{tab:cumplimiento} y~\ref{tab:rnf-cumplimiento}
del capítulo de resultados.  En síntesis, los cinco objetivos específicos
se cumplen y OE3 se supera (la \ac{API} creció a diez \emph{blueprints}
registrados con la incorporación de \texttt{patterns} en el Sprint~5);
los cinco requisitos no funcionales (seguridad, escalabilidad,
mantenibilidad, usabilidad y portabilidad) se atienden satisfactoriamente.
La seguridad, la escalabilidad y la usabilidad quedan como cumplimiento
parcial, pendientes de una auditoría de seguridad, de pruebas de carga y de
una evaluación de usabilidad respectivamente, junto con las pruebas
\emph{end-to-end} del \emph{frontend}.  Todo ello se recoge como trabajo
futuro.

\textbf{Valoración global del proyecto:} \dataqual es una plataforma
funcional de extremo a extremo cuyas fortalezas se detallan en la
sección~\ref{sec:fortalezas}.  Su madurez es la de un prototipo funcional,
no de una herramienta de producción, lo cual está alineado con el alcance
de un \ac{TFG}.

\textbf{Limitaciones encontradas:} las restricciones técnicas y de alcance
del sistema se analizan en la sección~\ref{sec:limitaciones}, y las líneas
de trabajo derivadas, en la sección~\ref{sec:trabajo-futuro}.

%% =========================================================================
\section{Fortalezas de la plataforma}
\label{sec:fortalezas}
%% =========================================================================

A partir de los resultados consolidados, las fortalezas de \dataqual se
pueden agrupar en cuatro ejes: una base arquitectónica extensible, un
conjunto de capacidades funcionales diferenciadoras, una integración
nativa de la protección de \acs{PII} y un modelo de despliegue y operación
robusto.  El desarrollo técnico de cada una se describe en los sprints
correspondientes del capítulo~\ref{chap:desarrollo}; aquí se discute su
valor desde la perspectiva del sistema acabado.

La principal fortaleza arquitectónica es el \textbf{patrón de
\emph{plugins} del motor de métricas}
(\texttt{BaseMetric}~+~\texttt{METRIC\_REGISTRY}, descrito en la
sección~\ref{sec:motor-metricas}).  Su valor quedó demostrado durante el
propio desarrollo, en el que el catálogo creció de tres métricas en el
Sprint~3 a siete al cierre del Sprint~5 sin necesidad de reescribir el
orquestador de evaluaciones.  La misma filosofía de extensibilidad se
traslada al \emph{frontend}, organizado en torno a más de cincuenta
componentes React reutilizables.

Sobre esa base, la plataforma incorpora cuatro \textbf{capacidades
funcionales diferenciadoras} frente a las herramientas analizadas en el
capítulo~\ref{chap:antecedentes}: el sistema de \emph{fingerprinting}
\ac{SHA}-256 (Sprint~5), que permite trazar la evolución de cada
\emph{issue} entre ejecuciones; los \emph{Quality Gates} con veredicto
\texttt{PASSED}/\texttt{WARNING}/\texttt{FAILED}, integrables en
\emph{pipelines} de \ac{MLOps}~\citep{sonarqube}; el módulo de perfilado
interactivo (Sprint~4), que concentra histogramas, diagramas de caja,
matrices de correlación y detección de \emph{outliers} en una sola
pantalla; y las \textbf{métricas específicas para \ac{ML}} (equilibrio
de clases y diversidad), ausentes en las herramientas generalistas.

La \textbf{protección nativa de \acs{PII}} es la fortaleza con mayor
valor diferencial en contextos regulados por el \ac{RGPD}.  Lejos de
tratarse de una capa filtro añadida \emph{a posteriori}, el
enmascaramiento es un contrato que toda métrica debe respetar
(Sprint~5): marca la diferencia entre «la herramienta puede ofuscar
\acs{PII} si se configura» y «la herramienta no puede revelar
\acs{PII} por construcción».

Desde el punto de vista operativo, \dataqual es \textbf{íntegramente
autoalojable mediante Docker Compose} (ocho servicios en una red
interna, sin dependencias de nube propietaria), incorpora un
\textbf{parseo robusto de \ac{CSV}} mediante matriz de \emph{fallback}
(Sprint~2), y orquesta el procesamiento asíncrono con
Celery~\citep{celery} complementado por un \emph{watchdog} y un
servicio híbrido de \emph{fallback} síncrono que garantizan la
completitud de las evaluaciones incluso en escenarios degradados.


%% =========================================================================
\section{Limitaciones y desafíos encontrados}
\label{sec:limitaciones}
%% =========================================================================

A pesar de los resultados alcanzados, la plataforma presenta limitaciones
que conviene reconocer explícitamente.  Estas se pueden agrupar en cinco
ámbitos: la escalabilidad y el rendimiento del motor de procesamiento,
la madurez de los servicios de plataforma de cara a entornos
corporativos, la cobertura de las pruebas automatizadas, el alcance de
la validación experimental y la frontera funcional del sistema en
relación con la corrección de datos y la integración con fuentes
externas.

El primer grupo afecta a la \textbf{escalabilidad y al rendimiento} bajo
cargas elevadas y constituye la restricción técnica más significativa de
la plataforma.  El motor de métricas opera sobre \emph{DataFrames} de
Pandas cargados íntegramente en memoria, lo que impone como cota
superior el tamaño de la RAM disponible y descarta procesar
\emph{datasets} de varios GB sin migrar a un motor distribuido como Dask
o Spark.  La capa de procesamiento, además, no se ha validado con
múltiples \emph{workers} de Celery en paralelo, por lo que la
escalabilidad horizontal que el diseño permite es todavía teórica.  El
\emph{frontend}, por su parte, obtiene el estado de las evaluaciones
mediante \emph{polling} periódico en lugar de notificaciones \emph{push}
por \emph{WebSocket}: aunque el patrón es funcionalmente correcto y
simplifica la arquitectura, introduce latencia perceptible y carga
innecesaria contra la \ac{API} cuando se ejecutan muchas evaluaciones
simultáneamente.  Por último, la detección de valores atípicos en el
módulo de perfilado se limita al método estadístico clásico del
\ac{IQR}; algoritmos basados en aprendizaje automático como
\emph{Isolation Forest} o \ac{DBSCAN}, que detectarían anomalías
multidimensionales más sutiles, quedan fuera del alcance de esta versión
(la decisión de tratar los \emph{outliers} como herramienta de
diagnóstico y no como dimensión de calidad se justifica en la
sección~\ref{sec:priorizacion-mvp}).

Un segundo grupo afecta a la \textbf{madurez de los servicios de
plataforma} de cara a entornos corporativos.  La autenticación se limita
a \ac{JWT} con usuario y contraseña: no se han implementado OAuth~2.0,
\emph{Single Sign-On} corporativo ni autenticación multifactor, por lo
que la plataforma no está preparada todavía para integrarse en entornos
empresariales con identidad federada.  Los resultados, además, no pueden
exportarse a formatos de informe formal como \ac{PDF} o Excel, lo que
dificulta el reporte a perfiles ajenos al equipo técnico y la
incorporación de los hallazgos a documentación de auditoría.

El tercer ámbito es la cobertura de \textbf{pruebas automatizadas}, que
aunque amplia en componentes individuales (motor de métricas,
\emph{Quality Score}, \emph{Quality Gates}, \emph{fingerprinting},
comparación con línea base, autorización \ac{IDOR} y versionado de
\emph{datasets}), no incluye pruebas \emph{end-to-end} que validen el
flujo completo desde la subida de un \emph{dataset} hasta la emisión del
veredicto del \emph{Quality Gate}.  Esta validación se realizó de forma
manual durante el desarrollo, pero su automatización con herramientas
como Playwright o Cypress queda como tarea pendiente y constituye la
mayor brecha en la estrategia de aseguramiento de calidad del propio
proyecto.

Finalmente, conviene reconocer también una limitación de carácter
metodológico relativa al \textbf{alcance de la validación
experimental}.  La verificación funcional del sistema se ha realizado
por el propio autor, complementada con un conjunto amplio de pruebas
automatizadas, pero no se ha llevado a cabo una validación con usuarios
externos del dominio de calidad de datos o de ingeniería de \ac{ML} que
permita evaluar empíricamente la usabilidad de la interfaz, la
adecuación del catálogo de métricas a casos de uso reales o el ajuste
de los umbrales por defecto de los \emph{Quality Gates}.  La
validación realizada acredita la corrección técnica del sistema, pero
no su idoneidad operativa para perfiles externos al desarrollo.

Un quinto ámbito es la \textbf{frontera funcional del sistema}.  Por un
lado, \dataqual diagnostica y clasifica los problemas de calidad, pero la
decisión y ejecución de las acciones correctoras (limpieza, imputación,
deduplicación) corresponden al usuario: la incorporación de sugerencias
de corrección automatizada queda como una de las líneas de trabajo
futuro identificadas en la sección~\ref{sec:trabajo-futuro}.  Por otro,
no se han desarrollado conectores específicos para bases de datos
(MySQL, BigQuery, Snowflake) ni para servicios de almacenamiento en la
nube (\ac{AWS}~S3, Azure~Blob, Google~Cloud~Storage); aunque la arquitectura
basada en MinIO es compatible con el protocolo S3, el usuario interactúa
actualmente con la plataforma a través de ficheros \ac{CSV} cargados
manualmente.

%% =========================================================================
\section{Líneas de trabajo futuro}
\label{sec:trabajo-futuro}
%% =========================================================================

A partir de las limitaciones identificadas en la
sección~\ref{sec:limitaciones} y de las oportunidades detectadas durante
el desarrollo, las líneas de evolución de \dataqual\ se articulan en
cuatro ejes: comunicación de resultados e integraciones salientes,
escalabilidad y rendimiento, integraciones empresariales, y
aseguramiento de calidad e inteligencia analítica.

En el eje de \textbf{comunicación de resultados e integraciones
salientes} se sitúa la exportación de informes a \ac{PDF} y Excel, que
facilitaría el reporte a \emph{stakeholders} no técnicos y la
incorporación de los hallazgos a documentación de auditoría.
Complementariamente, la incorporación de \emph{webhooks} permitiría
notificar a sistemas externos cuando una evaluación se complete o un
\emph{Quality Gate} falle, integrando \dataqual\ en flujos automatizados
de \ac{MLOps} sin necesidad de \emph{polling} desde el lado consumidor.

El eje de \textbf{escalabilidad y rendimiento} concentra tres mejoras
con cota técnica clara.  La sustitución del mecanismo actual de
\emph{polling} por notificaciones \emph{push} con \emph{WebSockets}
reduciría tanto la latencia perceptible como la carga innecesaria sobre
la \ac{API}.  El procesamiento por fragmentos (\emph{chunk-based})
levantaría la cota actual del tamaño de RAM disponible y abriría la
plataforma a \emph{datasets} de varios GB sin necesidad de migrar a un
motor distribuido como Dask o Spark.  Por último, las evaluaciones
programadas mediante Celery Beat (servicio ya disponible en la pila
tecnológica pero todavía sin interfaz de usuario) habilitarían la
monitorización continua de la calidad de datos sin intervención manual.

El tercer eje, \textbf{integraciones empresariales}, recoge las
limitaciones que separan a la plataforma de un escenario corporativo
real.  La incorporación de OAuth~2.0 y \emph{Single Sign-On} permitiría
integrar \dataqual\ en sistemas empresariales con identidad federada.
La integración directa con bases de datos relacionales y analíticas
(PostgreSQL, MySQL, BigQuery) como fuentes de datos eliminaría la
necesidad de exportar previamente los \emph{datasets} a fichero.  Y la
colaboración en equipo con roles y permisos granulares haría posible
que múltiples usuarios trabajaran sobre los mismos proyectos,
completando el modelo actual basado en un único propietario por
proyecto.

Por último, el eje de \textbf{aseguramiento de calidad e inteligencia
analítica} aborda tanto las pruebas del propio software como la
sofisticación analítica del producto.  La automatización de pruebas
\emph{end-to-end} con Cypress o Playwright cubriría los flujos
principales de la interfaz y cerraría la mayor brecha actual en la
estrategia de \ac{QA} del proyecto.  A ella se suman tres verificaciones
no funcionales pendientes: pruebas de carga que validen la escalabilidad
con grandes volúmenes de datos y usuarios concurrentes, una auditoría de
seguridad con pruebas de penetración, y una evaluación de usabilidad con
usuarios reales.  En el plano analítico, la
incorporación de algoritmos de detección de anomalías basados en
aprendizaje automático (\emph{Isolation Forest},
\emph{Autoencoders}) ampliaría el catálogo de métricas más allá del
método estadístico clásico \ac{IQR}, mientras que las sugerencias de
corrección automatizadas sobre los \emph{issues} detectados cerrarían
el ciclo de diagnóstico con propuestas concretas de acción.


%% =========================================================================
\section{Discusión}
\label{sec:discusion}
%% =========================================================================

Más allá de las fortalezas y limitaciones ya descritas, del trabajo se
desprenden dos conclusiones de fondo: la respuesta a una necesidad
práctica del ecosistema actual de calidad de datos y el valor de los
estándares internacionales como marco de rigor metodológico.

En primer lugar, el proyecto confirma que \textbf{la calidad de datos
para \ac{IA} requiere herramientas más accesibles} que los enfoques
\emph{ad hoc} basados en \emph{scripts} y cuadernos Jupyter.  Estas
aproximaciones, aunque flexibles, presentan barreras de estandarización,
historización y accesibilidad que dificultan su adopción sistemática en
organizaciones donde los responsables de la calidad no son perfiles
puramente técnicos.  \dataqual\ demuestra que es viable ofrecer una
interfaz web accesible sin renunciar a la potencia de una \ac{API}
\ac{REST}ful programática, conciliando ambos modos de uso bajo el mismo
sistema.

En segundo lugar, \textbf{los estándares internacionales han demostrado
ser un marco sólido para el diseño de métricas}: la alineación con
\ac{ISO}/\ac{IEC}~25012~\citep{iso25012} e
\ac{ISO}/\ac{IEC}~5259~\citep{iso5259} evitó definiciones arbitrarias y aportó
un vocabulario común que facilita la comparabilidad entre proyectos.


%% =========================================================================
\section{Lecciones aprendidas}
\label{sec:lecciones}
%% =========================================================================

El desarrollo del proyecto ha proporcionado lecciones que cabe agrupar
en dos planos: el metodológico, relativo a cómo se organiza el trabajo,
y el técnico, relativo a las decisiones de diseño que conviene tomar
pronto.

En el plano metodológico, la lección principal es que el \textbf{desarrollo
iterativo no es solo una metodología formal, sino una forma de pensar
que reduce el riesgo real de los proyectos}.  La organización en
iteraciones con entregas funcionales permitió detectar y corregir
problemas de diseño de forma temprana, evitando retrabajos costosos en
fases avanzadas.  Paralelamente, el desarrollo confirmó que \textbf{el
diseño de métricas de calidad de datos no es un ejercicio puramente
algorítmico}: requiere comprender los estándares internacionales, las
necesidades reales de los usuarios y los matices de cada dimensión de
calidad (qué significa exactamente «completitud» cuando una columna
tiene valores codificados como cadenas vacías, etiquetas
\texttt{N/A} o ceros con sentido distinto al de nulo).  Sin este trabajo
previo de modelado del dominio, las métricas resultantes habrían sido
superficiales o francamente engañosas.

En el plano técnico, dos decisiones tomadas pronto resultaron
especialmente rentables.  \textbf{La seguridad debe integrarse desde el
inicio del diseño}, no añadirse como capa posterior: decisiones como la
protección contra inyección en las reglas de consistencia lógica o el
enmascaramiento nativo de \ac{PII} fueron mucho más sencillas de
implementar al considerarlas desde las fases iniciales que si hubiera
sido necesario reconstruir el sistema sin ellas.  Por contraste, \textbf{la
gestión de la complejidad del \emph{frontend} requiere una disciplina
de modularización que conviene aplicar antes de que duela}: componentes
como el módulo de \emph{profiling} de datos crecieron significativamente
durante el desarrollo y, en retrospectiva, una mayor división temprana
en subcomponentes habría facilitado el mantenimiento en las iteraciones
finales.  La asimetría entre ambas lecciones resulta reveladora: lo que
se asume cuidar desde el principio (seguridad) sale bien, mientras que
lo que se posterga (refactor del \emph{frontend}) acaba doliendo.


%% =========================================================================
\section{Competencias adquiridas}
\label{sec:competencias}
%% =========================================================================

El desarrollo de este \ac{TFG} en la intensificación de Sistemas de
Información ha permitido adquirir y consolidar las siguientes competencias:

\begin{itemize}
  \item \textbf{Desarrollo \emph{full-stack}:} diseño e implementación de
    una aplicación web completa con separación clara entre \emph{frontend}
    (Next.js, React, TypeScript) y \emph{backend} (Python, Flask,
    SQLAlchemy).

  \item \textbf{Arquitecturas distribuidas y contenerización:} orquestación
    de ocho servicios Docker con Docker Compose, gestión de dependencias
    entre servicios, \emph{healthchecks} y volúmenes persistentes.

  \item \textbf{Procesamiento de datos tabulares:} manipulación de datasets
    en múltiples formatos con Pandas, cálculo de métricas estadísticas
    y gestión de casos extremos (nulos, tipos mixtos, codificaciones).

  \item \textbf{Calidad de datos e ingeniería de la calidad:} comprensión
    profunda de los estándares \ac{ISO}/\ac{IEC}~25012 e \ac{ISO}/\ac{IEC}~5259 y su
    aplicación práctica en el diseño de métricas.

  \item \textbf{Seguridad en aplicaciones web:} autenticación \ac{JWT},
    \emph{hashing} de contraseñas con \ac{PBKDF2}, limitación de tasa,
    cabeceras de seguridad y prevención de inyección de código.

  \item \textbf{Diseño de \ac{API}s \ac{REST}:} especificación y
    documentación de \emph{endpoints}, gestión de versiones, paginación y
    validación de entradas.

  \item \textbf{Gestión ágil con Scrum adaptado:} planificación de sprints,
    gestión del backlog con GitHub Issues, retrospectivas y adaptación del
    plan ante imprevistos.

  \item \textbf{Control de versiones avanzado:} estrategia de \emph{feature
    branches}, mensajes de \emph{commit} descriptivos y trazabilidad entre
    código e issues.

  \item \textbf{Estimación y planificación:} uso de \emph{story points},
    análisis de desviaciones y ajuste de estimaciones a lo largo del
    proyecto.

  \item \textbf{Documentación técnica:} redacción de una memoria de \ac{TFG}
    completa en \LaTeX, especificaciones de \ac{API} y catálogos de métricas.

  \item \textbf{Resolución autónoma de problemas:} diagnóstico y resolución
    de problemas de integración entre servicios Docker, incompatibilidades de
    formato de datos y gestión de errores en sistemas distribuidos.

  \item \textbf{Capacidad de síntesis y abstracción:} diseño de interfaces
    que permiten extender el sistema (motor de métricas, arquitectura de
    \emph{plugins}) sin conocer los detalles de las implementaciones concretas.

  \item \textbf{Trabajo autónomo bajo supervisión:} gestión del tiempo y la
    priorización en un proyecto individual de nueve meses con reuniones de
    seguimiento periódicas.
\end{itemize}


%% =========================================================================
\section{Valoración personal}
\label{sec:valoracion-personal}
%% =========================================================================

La elección del tema de este \ac{TFG} no fue casual: responde a mi interés
personal por el tratamiento y la gestión de los datos como disciplina, y a
mi voluntad de profundizar en el ámbito de la inteligencia artificial y la
ciencia de datos, campos en los que la calidad del dato de entrada resulta,
precisamente, el factor más determinante.  Llevar ese interés inicial al
terreno de una plataforma completa, y no quedarme en el plano teórico, es
lo que ha dado sentido al resto del trabajo.

Este \ac{TFG} ha supuesto un reto de una dimensión cualitativamente
diferente a los proyectos realizados durante el grado.  Por primera vez, he
tenido que diseñar y construir un sistema completo desde cero, tomando
decisiones de arquitectura con consecuencias a largo plazo, integrando por
primera vez en un sistema real tecnologías como Celery, MinIO y Redis
y manteniendo la coherencia de un proyecto de nueve meses en solitario.

Lo que más me ha sorprendido es la complejidad inherente a los detalles:
los ficheros \ac{CSV} del mundo real no se comportan como los de los tutoriales,
la configuración de Docker Compose con dependencias entre servicios exige
una comprensión profunda de los tiempos de arranque, y el diseño de un
\emph{fingerprint} determinista para issues requirió mucho más pensamiento
del esperado.  Estas dificultades, en retrospectiva, son precisamente las
más valiosas porque son las que ocurren en proyectos reales.

Me llevo especialmente la convicción de que una buena arquitectura inicial
(el modelo de datos pensado desde el principio, la arquitectura de
\emph{plugins} para el motor de métricas) ahorra enormemente el trabajo en
sprints posteriores, y que el desarrollo iterativo no es solo una metodología
formal sino una forma de pensar que reduce el riesgo real de los proyectos.

Estoy satisfecho con el resultado: \dataqual es una plataforma funcional,
segura y extensible que cubre un hueco real en el ecosistema de herramientas
de calidad de datos para \ac{IA}.  Queda trabajo por hacer, recogido en
la sección~\ref{sec:trabajo-futuro}, pero eso es precisamente lo que
convierte un \ac{TFG} en el punto de partida de algo mayor.


% Local Variables:
%  coding: utf-8
%  mode: latex
%  mode: flyspell
%  ispell-local-dictionary: "castellano8"
% End:
